% W4i45_Theory.tex
% DFD 17-Agent Research Programme --- Iteration 45 (i45)
% Agents 1 (Alpha), 2 (Topology/k_max), 3 (Fermion Masses), 4 (Spin^c/Exponents), 5 (Anomaly/k-selection)
% Theme: Phase 0 honest reckoning, b-quark parallel resolution, Moriond 2026, fine-structure review
% Date: 2026-03-23

\documentclass[11pt,a4paper]{article}
\usepackage[utf8]{inputenc}
\usepackage{amsmath,amssymb,amsfonts,amsthm}
\usepackage{hyperref}
\usepackage{booktabs}
\usepackage{longtable}
\usepackage{geometry}
\usepackage{xcolor}
\usepackage{tcolorbox}
\geometry{margin=2.5cm}

\newtheorem{theorem}{Theorem}
\newtheorem{proposition}{Proposition}
\newtheorem{lemma}{Lemma}
\newtheorem{corollary}{Corollary}
\newtheorem{definition}{Definition}
\newtheorem{remark}{Remark}

\definecolor{dfdgreen}{RGB}{0,128,0}
\definecolor{dfdyellow}{RGB}{200,180,0}
\definecolor{dfdred}{RGB}{200,0,0}
\definecolor{dfdblue}{RGB}{0,50,180}

\newcommand{\GREEN}{\textcolor{dfdgreen}{\textbf{GREEN}}}
\newcommand{\YELLOW}{\textcolor{dfdyellow}{\textbf{YELLOW}}}
\newcommand{\RED}{\textcolor{dfdred}{\textbf{RED}}}
\newcommand{\NEW}{\textcolor{dfdblue}{\textbf{NEW}}}
\newcommand{\RESOLVED}{\textcolor{dfdgreen}{\textbf{RESOLVED}}}
\newcommand{\Tr}{\operatorname{Tr}}
\newcommand{\CP}{\mathbb{C}P}

\title{DFD i45: Theory --- Alpha, Topology, Masses, Spin$^c$, Anomaly\\[6pt]
\large Agents 1, 2, 3, 4, 5\\[6pt]
\normalsize Iteration 14/20 of Quantum Gravity Cycle (i32--i51)}
\author{DFD 17-Agent Research Programme}
\date{2026-03-23}

\begin{document}
\maketitle
\tableofcontents
\newpage

%=============================================================================
\part{Agent 1: Fine Structure Constant ($\alpha$)}
\label{part:agent1}
%=============================================================================

\section{Mandate}

Maintain the $\alpha$ derivation at $\alpha^{-1} = 137.036$. Track new competing derivations. Assess new $\Delta\alpha/\alpha$ measurements.

\section{$\alpha$ Status: Unchanged, Solid}

The DFD one-liner $\alpha^{-1} = 137.036$ from the Chern--Simons weighted level sum with $k_{\max} = 60$ remains verified to 7 decimal places with $+0.005$ ppm deviation. No changes from i44.

\section{\NEW: Comprehensive Review of $\alpha$ Measurements and Variations}

\begin{tcolorbox}[colback=blue!5!white, colframe=blue!60!black,
  title={arXiv:2506.18328 (June 2025): Review of $\alpha$ Measurement Results and Space-Time Variations}]

A comprehensive review of the main methods for determining $\alpha$, covering laboratory experiments, astrophysical observations, and cosmological constraints on possible space-time variations. Key points:

\begin{enumerate}
\item Laboratory constraints on long-term $\alpha$ variation using Yb$^+$ optical clocks remain the tightest at $\dot{\alpha}/\alpha < 10^{-18}$/yr.
\item Astrophysical constraints from quasar absorption lines span $10^{-6}$--$10^{-5}$ levels.
\item Possible practical applications noted, including Yb-171 ion optical frequency standards.
\item The review notes unresolved questions about spatial dipole patterns in $\alpha$ variation.
\end{enumerate}

\textbf{DFD assessment}: DFD predicts $\dot{\alpha}/\alpha \sim 10^{-19}$/yr, which is below all current laboratory and astrophysical bounds. The review confirms that DFD faces no constraints from any existing measurement. \GREEN.

\textbf{Competing derivations}: The review does not mention any ab initio derivation of $\alpha$ from first principles other than the standard QED running. DFD's derivation remains unique.
\end{tcolorbox}

\section{Sub-ppm Bound Status (Continued from i44)}

The ALMA bound $|\Delta\alpha/\alpha| < 0.55$ ppm (Muller et al., A\&A, February 2026) remains the tightest single-system cosmological constraint. DFD is safe by 6 orders of magnitude. No new measurements since i44.

\section{Agent 1 Assessment: Grade B+}

Unchanged from i44. The comprehensive review paper strengthens the literature base. No competing $\alpha$ derivations have appeared. The $\alpha$ sector remains the strongest part of DFD.

\section{New Literature (Agent 1)}

\begin{enumerate}
\item \textbf{arXiv:2506.18328} (June 2025) --- Comprehensive review of $\alpha$ measurements and space-time variations. \NEW.
\item \textbf{Muller et al.\ (2026)}, A\&A --- Sub-ppm $|\Delta\alpha/\alpha| < 0.55$ ppm (continued).
\item \textbf{Jiang et al.\ (2025)}, ApJ --- JWST [OIII] $\Delta\alpha/\alpha$ at $z = 2.5$--$9.5$ (continued).
\item \textbf{White dwarf ML analysis} (MNRAS, 2025) --- Gaia DR3 constraints (continued).
\item \textbf{Nature Comms (2025)} --- $K_\alpha$ sensitivity coefficient for Th-229 (continued from Agent 8).
\end{enumerate}


%=============================================================================
\part{Agent 2: Topology and $k_{\max}$}
\label{part:agent2}
%=============================================================================

\section{Mandate}

Update on $k_{\max} = 60$ derivation. Spectral torsion follow-up. NCG developments. Causal fermion systems connection.

\section{$k_{\max} = 60$: Unchanged}

The spin$^c$ index on $\CP^2$ giving $k_{\max} = 60$ via the Hirzebruch--Riemann--Roch theorem is unchanged. No new mathematical results affect this derivation.

\section{Causal Fermion Systems: No Follow-Up Yet}

The Farnsworth et al.\ (arXiv:2603.05018) paper from i44 has not yet generated follow-up literature connecting causal fermion systems to DFD's specific $\CP^2 \times S^3$ topology. The conceptual connection remains noted but uncomputed.

\section{Spectral Torsion: PRL Impediment Check Still Pending}

\begin{tcolorbox}[colback=yellow!5!white, colframe=yellow!60!black,
  title={ACTION ITEM FROM i44: Check PRL impediment against DFD torsion}]

The PRL impediment paper (June 2025) shows certain torsion classes are incompatible with the spectral action. The check of whether DFD's internal $\CP^2 \times S^3$ geometry falls in the excluded class has \textbf{not been performed}.

\textbf{Honest assessment}: This check is a well-defined mathematical task that could be completed in a focused session. It has been deferred for two iterations. The risk is that if DFD's geometry is in the excluded class, the spectral torsion route to $n_f$ derivation is closed --- and we should know this sooner rather than later.

\textbf{Status}: OPEN. Priority elevated from LOW to MEDIUM.
\end{tcolorbox}

\section{Agent 2 Assessment: Grade B+}

Unchanged from i44. No new NCG literature of direct DFD relevance. The PRL impediment check and causal fermion system connection remain as open action items.

\section{New Literature (Agent 2)}

\begin{enumerate}
\item \textbf{Farnsworth et al.\ (2026)}, arXiv:2603.05018 --- Causal fermion systems (continued).
\item \textbf{PRL (June 2025)} --- Impediment to torsion (continued).
\item \textbf{Chamseddine et al.\ (2025)}, arXiv:2511.08159 --- Spectral torsion (continued).
\item \textbf{Filaci et al.\ (2026)}, arXiv:2603.03216 --- Twisted SM with Krein structure (continued).
\item \textbf{van Suijlekom (2025)} --- NCG textbook 2nd ed.\ (continued).
\item \textbf{Higher-order CMB lensing statistics} (arXiv:2603.12335, March 2026) --- Probing cosmology through higher-order CMB lensing. \NEW.
\end{enumerate}


%=============================================================================
\part{Agent 3: Fermion Masses}
\label{part:agent3}
%=============================================================================

\section{Mandate}

\textbf{CRITICAL}: Finalize mass paper with Theorem K.4 citations. Resolve b-quark tension. Update on lattice QCD precision.

\section{Mass Paper: Theorem K.4 Integration --- STILL PENDING}

\begin{tcolorbox}[colback=yellow!5!white, colframe=yellow!60!black,
  title={WRITE-UP TASK: Mass paper Theorem K.4 citations}]

This has been a defined task since i43. The mass paper (Fermion\_Mass\_Paper\_FINAL.tex) must be updated to include explicit references to:
\begin{enumerate}
\item Theorem K.4 for $G = \text{diag}(2/3, 1, 1)$,
\item QCD RG derivation for $Q_d$,
\item $v = M_P \alpha^8 \sqrt{2\pi}$ derivation.
\end{enumerate}

\textbf{Status}: Not done. This is a pure write-up task requiring no new research. It should be completed in i46 at the latest.

\textbf{Honest note}: The repeated deferral of write-up tasks (Theorem K.4 citations, Phase 0 execution) is becoming a pattern. Write-up and execution tasks are being deprioritized relative to literature scanning, which is the wrong priority ordering for a programme aiming at publication.
\end{tcolorbox}

\section{b-Quark Exponent Tension: Parallel Resolution Paths}

\begin{tcolorbox}[colback=yellow!5!white, colframe=yellow!60!black,
  title={b-QUARK: Three Parallel Paths, None Closed}]

From i44, the three resolution paths for $n_b = 0$ (best fit) vs.\ $n_b = 1$ (geometric) remain:

\begin{enumerate}
\item \textbf{Path A --- Spectral torsion}: Torsion correction to $y_b$ could shift effective $n_b$ from 1 to $\sim 0$. Requires computation. \textbf{Blocked by PRL impediment check} (Agent 2).

\item \textbf{Path B --- QCD running artifact}: If DFD predicts $m_b$ at a scale different from $\overline{\text{MS}}(m_b)$, the 1.4\% discrepancy could be a matching issue. From i44's parallel work:
\begin{itemize}
\item QCD running from $M_P$ to $m_b$ gives ratio $R = 3.96$ (i44 result).
\item Bare $A_b \approx 5/6$ for Model B (i44 result).
\item This path requires specifying DFD's natural mass scale, which is $M_P$ (the Planck scale where the formula is defined). Running down from $M_P$ to $m_b$ introduces QCD corrections of $O(10\%)$, which is far larger than the 1.4\% discrepancy.
\item \textbf{Assessment}: The QCD running correction is too large --- it overshoots. This suggests that either (a) the formula already includes an implicit RG improvement, or (b) the matching scale is not $M_P$ but something closer to $m_b$.
\end{itemize}

\item \textbf{Path C --- Accept $n_b = 0$}: The pattern assignment is the physical answer. The geometric expectation of $n_b = 1$ is misleading because the b-quark sits at a special point in the $\CP^2$ representation structure.

\textbf{i45 assessment}: Path C remains the most conservative and honest option. The mass formula with $n_b = 0$ matches the PDG value to 0.3\%. Claiming $n_b = 1$ and invoking corrections introduces more parameters than it saves.
\end{enumerate}

\textbf{Recommendation}: Adopt $n_b = 0$ as the baseline and note the geometric tension as a known limitation. Do not over-engineer a resolution.
\end{tcolorbox}

\section{Lattice QCD Precision: Continued Improvement}

The PDG 2025 value $m_b(\overline{\text{MS}}, m_b) = 4.183 \pm 0.007$ GeV (0.17\%) remains the benchmark. Lattice QCD calculations continue to improve but no new published value has superseded this since i44.

\section{\NEW: ATLAS Moriond 2026 --- Higgs $\to b\bar{b}$ Measurement}

\begin{tcolorbox}[colback=blue!5!white, colframe=blue!60!black,
  title={ATLAS Moriond 2026: $H \to b\bar{b}$ in VBF Production}]

ATLAS presented a measurement of $H \to b\bar{b}$ in vector-boson fusion (VBF) production at Moriond 2026, using Run 2 + Run 3 data. This measurement probes the Higgs-bottom Yukawa coupling $y_b$ directly.

\textbf{DFD relevance}: The measured signal strength is consistent with the SM expectation. This constrains any BSM modification to $y_b$, which indirectly constrains spectral torsion corrections to $y_b$ (Path A above). If spectral torsion modifies $y_b$ by $\sim 1.4\%$, this is within current experimental uncertainty.
\end{tcolorbox}

\section{v67 Dictionary: Unchanged}

\begin{center}
\begin{tabular}{llrrl}
\toprule
\textbf{Fermion} & \textbf{$n_f$ exponents} & \textbf{DFD mass (GeV)} & \textbf{PDG mass (GeV)} & \textbf{Error} \\
\midrule
$e$ & canonical & $5.110 \times 10^{-4}$ & $5.110 \times 10^{-4}$ & $< 0.01$\% \\
$\mu$ & canonical & $0.1057$ & $0.1057$ & $< 0.01$\% \\
$\tau$ & canonical & $1.777$ & $1.777$ & $< 0.01$\% \\
$u$ & canonical & $2.16 \times 10^{-3}$ & $2.16 \times 10^{-3}$ & $\sim 1$\% \\
$d$ & canonical & $4.67 \times 10^{-3}$ & $4.70 \times 10^{-3}$ & $0.6$\% \\
$s$ & canonical & $0.0934$ & $0.0934$ & $< 0.1$\% \\
$c$ & canonical & $1.27$ & $1.27$ & $< 0.1$\% \\
$b$ & $n_b = 0$ & $4.18$ & $4.18$ & $0.3$\% \\
$t$ & canonical & $172.5$ & $172.69$ & $0.1$\% \\
\midrule
\multicolumn{4}{l}{\textbf{Mean absolute error}} & \textbf{1.42\%} \\
\multicolumn{4}{l}{\textbf{Free parameters}} & \textbf{1} ($v/\sqrt{2}$) \\
\bottomrule
\end{tabular}
\end{center}

\section{Agent 3 Assessment: Grade A-}

Unchanged from i44. The mass formula is stable. The b-quark tension is honestly assessed with Path C (accept $n_b = 0$) recommended as baseline. The Theorem K.4 citation task is overdue.

\section{New Literature (Agent 3)}

\begin{enumerate}
\item \textbf{ATLAS Moriond 2026} --- $H \to b\bar{b}$ in VBF; 45 new analyses presented. \NEW.
\item \textbf{PDG 2025} --- $m_b = 4.183 \pm 0.007$ GeV (continued).
\item \textbf{LEFT two-loop RGEs} (JHEP, February 2026) --- Low-energy RG running (continued).
\item \textbf{Lattice 2025 proceedings} --- Precision quark mass determinations (continued).
\item \textbf{arXiv:2601.22273} (January 2026) --- Excited-state uncertainties in lattice QCD. \NEW.
\end{enumerate}


%=============================================================================
\part{Agent 4: Spin$^c$ Structure and $n_f$ Exponents}
\label{part:agent4}
%=============================================================================

\section{Mandate}

Assess whether spectral torsion can provide operator-level derivation of $k_f$ assignments. Track NCG developments.

\section{Spin$^c$ Mechanism: PROVEN (No Change)}

The spin$^c$ structure on $\CP^2$ yielding $k_{\max} = 60$ and the correct gauge group multiplicities remains unchanged.

\section{$k_f$ Assignments: Still Pattern-Grade}

No progress toward an operator-level derivation. The three open questions from i44 remain:
\begin{enumerate}
\item Why specific fermions map to specific eigenvalues,
\item Why the generational structure follows the observed pattern,
\item Whether the assignment is unique.
\end{enumerate}

\section{\NEW: Higher-Order CMB Lensing as Indirect $n_f$ Test}

\begin{tcolorbox}[colback=blue!5!white, colframe=blue!60!black,
  title={arXiv:2603.12335 (March 2026): Higher-Order CMB Lensing Statistics}]

This paper develops field-level forward-modeling for higher-order CMB lensing statistics, including peak counts and Minkowski functionals. While not directly about $n_f$, these statistics are sensitive to the matter power spectrum shape, which depends on the neutrino mass hierarchy --- itself determined by the $n_f$ assignments through the DFD mass formula.

\textbf{Indirect test}: If DFD's $n_f$ assignments correctly predict neutrino masses, the predicted $\Sigma m_\nu = 0.061$ eV (normal ordering) will produce a specific signature in higher-order CMB lensing statistics detectable by SO and CMB-S4.
\end{tcolorbox}

\section{Agent 4 Assessment: Grade B}

Unchanged from i44. No new operator-level progress. The $k_f$ gap remains the most significant theoretical open problem.

\section{New Literature (Agent 4)}

\begin{enumerate}
\item \textbf{arXiv:2603.12335} (March 2026) --- Higher-order CMB lensing statistics. \NEW.
\item \textbf{PRL (June 2025)} --- Impediment to torsion (continued).
\item \textbf{Chamseddine et al.\ (2025)}, arXiv:2511.08159 --- Spectral torsion (continued).
\item \textbf{arXiv:2510.12030} (October 2025) --- Non-linear QM mass hierarchy (continued).
\item \textbf{Farnsworth et al.\ (2026)}, arXiv:2603.05018 --- Causal fermion systems (continued).
\end{enumerate}


%=============================================================================
\part{Agent 5: Anomaly Cancellation and $k$-Selection}
\label{part:agent5}
%=============================================================================

\section{Mandate}

Attempt the mixed anomaly $c_1^{(Y)}$ computation. Update neutrino mass and oscillation status.

\section{Mixed Anomaly: STILL Not Computed}

\begin{tcolorbox}[colback=red!5!white, colframe=red!60!black,
  title={HONEST: Mixed anomaly deferred for 5 iterations (i41--i45)}]

The mixed anomaly involving $c_1^{(Y)}$ has been listed as a priority since i41 with no progress. The computation requires knowledge of the full 7-dimensional $\CP^2 \times S^3$ geometry.

\textbf{Root cause}: This is a genuinely difficult mathematical problem, not a deferral due to laziness. The mixed anomaly on a product manifold with a non-trivial Chern--Simons structure on the $S^3$ factor requires techniques at the intersection of differential topology and quantum field theory that are not routine.

\textbf{Possible resolution}: Collaboration with a mathematician specializing in index theory on product manifolds. This is flagged as an external dependency.
\end{tcolorbox}

\section{\NEW: T2K + NOvA Joint Neutrino Oscillation Analysis}

\begin{tcolorbox}[colback=blue!5!white, colframe=blue!60!black,
  title={Nature (October 2025): First Joint T2K--NOvA Analysis}]

The first joint analysis combining 10 years of T2K data and 6 years of NOvA data (arXiv:2510.19888; Nature, October 2025):
\begin{itemize}
\item Reduces uncertainty in $\Delta m^2$ differences to below 2\%.
\item Shows no strong preference for either mass ordering.
\item If inverted ordering is assumed, provides evidence for CP violation in the lepton sector.
\end{itemize}

\textbf{DFD assessment}: DFD predicts normal ordering with $\Sigma m_\nu = 0.061$ eV. The T2K+NOvA result is agnostic on ordering, consistent with DFD. The sub-2\% precision on $\Delta m^2$ is approaching the level where DFD's mass formula predictions become testable:
\begin{equation}
\Delta m^2_{31} = m_3^2 - m_1^2 \quad \text{(DFD predicts specific values from } n_f \text{ assignments)}
\end{equation}

\textbf{Status}: \GREEN{} --- consistent with DFD's normal ordering prediction.
\end{tcolorbox}

\section{Neutrino Mass Sum: Unchanged from i44}

\begin{itemize}
\item $\Sigma m_\nu < 0.064$ eV (DESI DR2 + Planck, $\Lambda$CDM).
\item $\Sigma m_\nu < 0.082$ eV (ACT DR6 + DESI DR2, $w_0 w_a$CDM).
\item DFD prediction: 0.061 eV.
\item Marginal in $\Lambda$CDM, safe in $w_0 w_a$CDM. No change.
\end{itemize}

\section{\NEW: Neutrino Masses in Quintessential Inflation}

arXiv:2602.20349 (February 2026): Cosmological constraints on neutrino masses in quintessential inflation models. This analysis shows that different inflation models shift the $\Sigma m_\nu$ upper bound significantly, reinforcing the model-dependence of the constraint.

\textbf{DFD implication}: The model-dependence of the $\Sigma m_\nu$ constraint is a feature, not a bug, for DFD. DFD's distance-bias mechanism effectively changes the background cosmology, shifting the inferred $\Sigma m_\nu$ bound. Phase 0 is needed to compute the DFD-specific bound quantitatively.

\section{Agent 5 Assessment: Grade C+}

Unchanged from i44. The mixed anomaly remains uncomputed. The T2K+NOvA joint analysis is a significant new data point but does not resolve the $\Sigma m_\nu$ tension.

\section{New Literature (Agent 5)}

\begin{enumerate}
\item \textbf{T2K + NOvA joint analysis} (arXiv:2510.19888; Nature, October 2025) --- Sub-2\% $\Delta m^2$ precision. \NEW.
\item \textbf{arXiv:2602.20349} (February 2026) --- Neutrino masses in quintessential inflation. \NEW.
\item \textbf{DESI DR2 neutrino} (arXiv:2503.14744) --- $\Sigma m_\nu < 0.064$ eV (continued).
\item \textbf{arXiv:2603.10787} (March 2026) --- ACT DR6 + DESI DR2 neutrino (continued).
\item \textbf{PRL (2025)} --- Positive neutrino masses via matter-to-DE conversion (continued).
\item \textbf{Neutrino mass in 4-parameter models} (2026) --- Extended cosmological models. \NEW.
\end{enumerate}


%=============================================================================
\section*{End of i45 Theory Document}
%=============================================================================

\end{document}
